\documentclass[11pt,a4paper]{article}
\usepackage[utf8]{inputenc}
\usepackage[T1]{fontenc}
\usepackage[brazil]{babel}
\usepackage{geometry}
\usepackage{graphicx}
\usepackage{booktabs}
\usepackage{longtable}
\usepackage{array}
\usepackage{float}
\usepackage{xcolor}
\usepackage{listings}
\usepackage{pdflscape}
\usepackage{placeins}
\usepackage[hidelinks]{hyperref}
\geometry{margin=2.1cm}
\setlength{\parskip}{0.5em}
\setlength{\parindent}{0pt}
\lstset{
  basicstyle=\ttfamily\footnotesize,
  breaklines=true,
  columns=fullflexible,
  frame=single,
  showstringspaces=false,
  keywordstyle=\color{blue},
  commentstyle=\color{gray}
}
\title{Guia de métricas do benchmark Druid}
\author{CCO DataOps}
\date{30 de junho de 2026}
\begin{document}
\maketitle


\section{Objetivo do guia}

Este guia define as métricas usadas nos relatórios de capacidade. Ele existe para evitar leituras equivocadas, como comparar QPS ignorando erro ou interpretar p95 baixo em um cenário que perdeu conexões.

\section{Glossário operacional}

{
\footnotesize
\setlength{\tabcolsep}{4pt}
\renewcommand{\arraystretch}{1.15}
\begin{longtable}{p{0.14\textwidth}p{0.28\textwidth}p{0.28\textwidth}p{0.22\textwidth}}
\caption{Glossario de metricas usadas na matriz.}\\
\toprule
Metrica & Definicao & Como ler & Cuidado \\
\midrule
\endfirsthead
\toprule
Metrica & Definicao & Como ler & Cuidado \\
\midrule
\endhead
QPS & Consultas finalizadas por segundo dentro do bloco medido. & Maior QPS indica maior vazao, mas so e ganho real se erro e p95 continuarem aceitaveis. & QPS pode parecer bom quando muitas consultas falham rapidamente. \\
p50 & Mediana da latencia: metade das requisicoes foi mais rapida que esse valor. & Mostra a experiencia tipica de consulta. & Nao revela cauda longa nem saturacao. \\
p95 & Latencia abaixo da qual 95\% das requisicoes ficaram. & E a principal metrica de estabilidade percebida em dashboards. & Deve ser lida junto com taxa de erro. \\
p99 & Latencia de cauda extrema: 99\% das requisicoes ficaram abaixo dela. & Indica o custo dos piores casos ainda bem-sucedidos. & Com poucas amostras, oscila bastante; ainda assim revela risco operacional. \\
Taxa de erro & Percentual de requisicoes que terminaram com excecao, timeout ou erro HTTP. & Qualquer valor acima de zero em benchmark pequeno merece investigacao. & Nao compare latencia entre cenarios ignorando erros. \\
Saturacao & Classificacao derivada de erro ou estouro dos limites configurados no runner. & Marca quando o cenario deixou de sustentar a carga de forma confiavel. & Uma rodada saturada nao deve ser vendida como ganho de performance. \\
Cold-start & Execucao com cache de consulta desabilitado para aproximar leitura fria. & Mostra o custo de buscar e processar dados sem reaproveitamento. & Nao representa exatamente reinicio fisico de containers. \\
Warm-start & Execucao com cache habilitado e populado para medir reaproveitamento. & Mostra ganho quando dashboards repetem padroes de consulta. & Cache nao corrige falta de CPU/memoria sob alta concorrencia. \\
CPU Docker & Percentual de CPU usado pelo container durante a matriz. & Picos acima de 100\% indicam uso de mais de um core logico. & Compare com limites de CPU do perfil, nao com a maquina inteira apenas. \\
Memoria Docker & Percentual de memoria usada dentro do limite do container. & Queda percentual pode ocorrer porque o limite aumentou, nao porque a carga diminuiu. & Leia junto com CPU e erros do Historical. \\
Loadstatus & Percentual de segmentos carregados e disponiveis no Druid. & 100\% indica que os dados estavam disponiveis para consulta. & Loadstatus bom nao garante performance; apenas disponibilidade. \\
\bottomrule
\end{longtable}
}


\section{Como combinar métricas}

\subsection{QPS não deve ser lido sozinho}

QPS mede vazão, mas não qualidade. Se duas configurações entregam QPS parecido, a melhor é a que mantém menor erro e menor p95. Se uma configuração aumenta QPS mas passa a falhar, ela atingiu saturação operacional.

\subsection{p95 e p99 explicam experiência de painel}

A média pode parecer confortável quando poucas consultas são muito lentas. p95 e p99 mostram a cauda, que é o que bloqueia dashboards compostos: a tela espera o card mais lento.

\subsection{Erro tem prioridade interpretativa}

Quando há erro, a latência dos sucessos não representa a carga inteira. Por isso Baseline e Balanced em concorrência 4 devem ser classificados como saturados mesmo quando algumas consultas retornaram rápido.

\subsection{Cold e warm start respondem perguntas diferentes}

Cold-start aproxima o custo de leitura sem reaproveitamento. Warm-start aproxima repetição de painéis e filtros. Um bom perfil deve se beneficiar de warm-start, mas continuar confiável quando a consulta fria aparecer.

\subsection{CPU e memória mostram o mecanismo da falha}

As métricas Docker não substituem as métricas SQL; elas explicam por que a curva se comportou daquele modo. Na matriz, Broker e Historical sobem CPU no Query Heavy e deixam de apresentar reset de canal.

\section{Critério de leitura adotado}

\begin{enumerate}

\item Verificar se loadstatus estava completo.

\item Verificar taxa de erro.

\item Comparar QPS apenas entre cenários sem erro equivalente.

\item Comparar p95/p99 para entender cauda operacional.

\item Usar Docker stats para explicar saturação.

\end{enumerate}


\end{document}
